\documentclass{training}
\setlang{french}
% Filigrane PAR MODULE : \uxfond (PNG A4 dilue 30 %). Cadre noir vectoriel commun.
% Proportions original 744 px : marge 20/744, trait 1/744. Exterieur BLANC.
% Cadre via eso-pic +\put (coordonnees absolues) — pas de remember picture (decalage).
\usepackage{eso-pic}

\providecommand{\uxfond}{../../resources/fond/norme_fond_paliers_30.png}

\newlength{\uxmargin}
\newlength{\uxthick}
\newlength{\uxinnerw}
\newlength{\uxinnerh}
\setlength{\uxmargin}{0.02688\paperwidth}% 20/744
\setlength{\uxthick}{0.001344\paperwidth}% 1/744
\setlength{\uxinnerw}{\dimexpr\paperwidth-2\uxmargin-2\uxthick\relax}
\setlength{\uxinnerh}{\dimexpr\paperheight-2\uxmargin-2\uxthick\relax}

% Contenu strictement a l'interieur du trait (+ jour confortable).
\geometry{%
  includeheadfoot,
  left=22mm,
  right=22mm,
  top=20mm,
  bottom=18mm,
  headheight=30pt,
  headsep=5mm,
  footskip=12mm%
}

\AddToShipoutPictureBG{%
  \AtPageLowerLeft{%
    % Fond blanc page entiere
    \put(0,0){\color{white}\rule{\paperwidth}{\paperheight}}%
    % Filigrane module a l'interieur du cadre
    \put(\LenToUnit{\dimexpr\uxmargin+\uxthick\relax},\LenToUnit{\dimexpr\uxmargin+\uxthick\relax}){%
      \includegraphics[width=\uxinnerw,height=\uxinnerh]{\uxfond}}%
    % Cadre noir : bas, haut, gauche, droite
    \put(\LenToUnit{\uxmargin},\LenToUnit{\uxmargin}){%
      \color{black}\rule{\dimexpr\paperwidth-2\uxmargin\relax}{\uxthick}}%
    \put(\LenToUnit{\uxmargin},\LenToUnit{\dimexpr\paperheight-\uxmargin-\uxthick\relax}){%
      \color{black}\rule{\dimexpr\paperwidth-2\uxmargin\relax}{\uxthick}}%
    \put(\LenToUnit{\uxmargin},\LenToUnit{\uxmargin}){%
      \color{black}\rule{\uxthick}{\dimexpr\paperheight-2\uxmargin\relax}}%
    \put(\LenToUnit{\dimexpr\paperwidth-\uxmargin-\uxthick\relax},\LenToUnit{\uxmargin}){%
      \color{black}\rule{\uxthick}{\dimexpr\paperheight-2\uxmargin\relax}}%
  }%
}

\makeatletter
\fancyhf{}
\fancyheadoffset{0pt}
\renewcommand{\headrulewidth}{0pt}
\lhead{\parbox[t]{0.56\textwidth}{\RaggedRight\footnotesize\sloppy\@title}}
\rhead{\parbox[t]{0.36\textwidth}{\RaggedLeft\footnotesize\sloppy\@subtitle}}
\cfoot{\thepage}
% Coupures : underscore et slash dans les noms de fichiers / consignes
\DeclareRobustCommand{\_}{\textunderscore\allowbreak}
\makeatother

% Evite les depassements de longues lignes (Github, Name_FirstName_...)
\sloppy
\emergencystretch=3em
\hbadness=10000

% Sommaire type 42 : chapitres seulement
\setcounter{tocdepth}{0}
\renewcommand{\cftchappresnum}{}
\renewcommand{\cftchapaftersnum}{}
\setlength{\cftchapnumwidth}{2.0em}
\setlength{\cftbeforechapskip}{0.35em}
\cftsetindents{chapter}{0em}{2.0em}


\begin{document}

\schoollogo{../../resources/jeunecrack.png}
\title{Règles}
\subtitle{Piscine UX/UI --- UI}
\member{Piscine UX/UI}{https://github.com/01-edu/UX-UI}
\summary{Grilles, états, accessibilité, et les petits contrôles qu'on réutilise.}
\version{1.0}
\prerequisite{Le module précédent de cette piscine, ou aucun si c'est le premier.}
\maketitle

\newpage
\tableofcontents
\newpage


\chapter{Instructions}

L'avant-propos n'est pas un exercice. Il ne sera pas noté. Les exercices commencent au chapitre suivant, du contrat le plus simple au plus dur. Un exercice plus tardif n'est pas pris en compte si un plus tôt ne tient pas.

On ne regarde que ce qui est dans le répertoire de rendu de l'exercice. N'ajoutez pas un fichier que l'exercice ne nomme pas. Le nom d'un livrable suit le motif imprimé dans l'exercice : \texttt{Name\_FirstName\_DeliverableName\_Date\_VersionNumber}.

Le geste ne vous est pas donné. Les exemples d'un exercice portent parfois une contrainte que la phrase ne répète pas. Lisez-les. Vos autres références sont le manuel de l'outil choisi, la personne assise à côté, et les liens imprimés dans l'exercice.

L'ordre est celui de la quête. Sauter une étape parce qu'un écran plus tard semble plus intéressant produit un écran que personne ne sait expliquer.

\subsection*{Des images}

La marque en couverture est le logo du projet. L'image de l'avant-propos est l'illustration de ce module. Le fond pâle derrière le texte est le filigrane de page. Aucun n'est un indice.

\subsection*{Des machines qui écrivent}

Vous pouvez vous servir d'un outil pour vérifier un mot, comme d'un dictionnaire. Vous ne pouvez pas remplacer un entretien, un test ou un fichier par un texte généré et le signer. Si vous ne pouvez pas montrer la personne, l'écran ou la ligne que vous avancez, l'exercice n'est pas fait.


\chapter{Avant-propos}

Josef Müller-Brockmann a écrit \emph{Grid Systems in Graphic Design} pour que l'alignement soit une décision et non une habitude. \emph{The Vignelli Canon} de Massimo Vignelli est la version courte du même entêtement. Vous tracerez des lignes sur les écrans des autres. Les livres ne les tracent pas.

\begin{center}
\includegraphics[width=0.78\textwidth,height=0.32\textheight,keepaspectratio]{../../resources/intros/UI-02-rules.jpg}\\[0.4em]
{\footnotesize Image : AlanKitching / CC BY 4.0}
\end{center}

À regarder. Les vidéos sont un contexte. Elles ne sont pas le livrable.
\begin{itemize}
\item UI - Introduction \& Quest n°1, 01Edu \url{https://www.youtube.com/watch?v=UbE6WMeuvQg&list=PLHyAJ_GrRtf-_aBft2jQL4-QAunn2bLvu}
\end{itemize}

\chapter{Exercice 00 : UI III - Ex 1 - Grilles}
\exnumber{00}
\extitle{UI III - Ex 1 - Grilles}
\exfiles{ex00/fichier Figma, nomme Name\_FirstName\_DeliverableName\_Date\_VersionNumber}
\exauthorize{Ceux nommés dans l'exercice}
\makeheaderfiles

\textbf{Contexte :}

Bienvenue dans la 3e quête de UI.

Dans cette série d'exercices, on vous demandera d'ouvrir les yeux et de collecter des éléments graphiques sur divers sites web. Vous verrez des grilles à travers les écrans, des boutons multi-états, comment rendre un design accessible pour tous, des fils d'Ariane, des boutons toggle, des boutons radio, un calendrier, un time picker, avec 2 objectifs principaux :

\begin{itemize}
\item Enrichir votre culture générale graphique
\item Détecter la cohérence, qui aide les utilisateurs à comprendre un produit numérique dès le premier regard.
\end{itemize}

\textbf{Instructions :}

Ouvrez ces 5 images sur Figma et tracez des lignes et des grilles chaque fois que vous trouvez des éléments alignés.

\textbf{Livrables :}

\begin{itemize}
\item Un fichier Figma comprenant les écrans avec des lignes par-dessus pour montrer les grilles. Une page par écran ou frame.
\item Déposez le fichier Figma sur Github, nommé ainsi : « Name\_FirstName\_DeliverableName\_Date\_VersionNumber ». Exemple : « Doe\_John\_MidFiPrototype\_05242024\_V1 ».
\item Ajoutez un titre dans le fichier.
\item Soignez l'aspect global du fichier. Il doit être clair, simple et lisible. Vous pouvez regarder canva.com, sans surcharger le dessin de détails.
\end{itemize}

\textbf{Ressources :}

\begin{itemize}
\item Tutoriel vidéo sur Figma \url{https://www.youtube.com/watch?v=FTFaQWZBqQ8}
\item Débuter en design graphique \url{https://edu.gcfglobal.org/en/beginning-graphic-design/}
\item Gareth graphic studio \url{https://www.youtube.com/c/GarethDavidStudio}
\item Exemple du 1er écran \url{https://www.figma.com/file/KOzd8yZuJvGLmZLuXjBgDg/UI-III-Ex-1.1}
\end{itemize}


\chapter{Exercice 01 : UI III - Ex 2 - États multiples}
\exnumber{01}
\extitle{UI III - Ex 2 - États multiples}
\exfiles{ex01/fichier Figma, nomme Name\_FirstName\_DeliverableName\_Date\_VersionNumber}
\exauthorize{Those named in the exercise}
\makeheaderfiles

\textbf{Instructions}:

Go to Sony.com \url{http://sony.com} and observe the buttons and call to actions. When you hover or click on them, they change!

Make screenshots of the different states of buttons, menu titles, filters etc. Then, gather and organize them.

Find at least five combinations of multi state elements.

\textbf{Livrables :}

\begin{itemize}
\item A Figma file with 5 combinations of multi state elements.
\item Déposez le fichier Figma sur Github, nommé ainsi : « Name\_FirstName\_DeliverableName\_Date\_VersionNumber ». Exemple : « Doe\_John\_MidFiPrototype\_05242024\_V1 ».
\item Ajoutez un titre dans le fichier.
\item Soignez l'aspect global du fichier. Il doit être clair, simple et lisible. Vous pouvez regarder canva.com, sans surcharger le dessin de détails.
\end{itemize}

\textbf{Ressources :}

\begin{itemize}
\item Video tutorial on Figma \url{https://www.youtube.com/watch?v=FTFaQWZBqQ8}
\item Beginning graphic design \url{https://edu.gcfglobal.org/en/beginning-graphic-design/}
\item Gareth graphic studio \url{https://www.youtube.com/c/GarethDavidStudio}
\item Example here with Nintendo.com \url{https://www.figma.com/file/ahSnWRZeKBO62oJDiXltxY/UI-III---Ex-2}
\end{itemize}

/!\textbackslash{} The screenshots date back to nov 2021. They may differ from the current version of Nintendo.com \url{http://Nintendo.com} but the rules are the same


\chapter{Exercice 02 : UI III - Ex 3 - Accessibilité (site)}
\exnumber{02}
\extitle{UI III - Ex 3 - Accessibilité (site)}
\exfiles{ex02/fichier Figma, nomme Name\_FirstName\_DeliverableName\_Date\_VersionNumber}
\exauthorize{Ceux nommés dans l'exercice}
\makeheaderfiles

\textbf{Instructions :}

Choisissez un site web que vous trouvez particulièrement accessible. Trouvez au moins 3 exemples concrets de mesures d'accessibilité et expliquez comment elles facilitent la lecture ou la perception du contenu visuel.

\textbf{Livrables :}

\begin{itemize}
\item Un fichier Figma de 3 pages avec 3 éléments différents issus d'un site (ils peuvent venir de 3 sites différents) qui sont accessibles. Les trois mesures doivent être différentes.
\item Déposez le fichier Figma sur Github, nommé ainsi : « Name\_FirstName\_DeliverableName\_Date\_VersionNumber ». Exemple : « Doe\_John\_MidFiPrototype\_05242024\_V1 ».
\item Ajoutez un titre dans le fichier.
\item Soignez l'aspect global du fichier. Il doit être clair, simple et lisible. Vous pouvez regarder canva.com, sans surcharger le dessin de détails.
\end{itemize}

\textbf{Ressources :}

\begin{itemize}
\item Ce qu'est l'accessibilité et pourquoi elle est si importante \url{https://uxdesign.cc/what-accessibility-is-and-why-its-so-important-9c56e033ff26}
\item Web Content Accessibility Guidelines (WCAG) \url{https://www.w3.org/WAI/standards-guidelines/wcag/}
\item Tests d'accessibilité  ---  W3C Wik \url{https://www.w3.org/wiki/Accessibility_testing}i
\item Mythe UX : l'accessibilité est chère et difficile \url{https://uxmyths.com/post/654091803/myth-5-accessibility-is-expensive-and-difficult}
\item Salesforce UX - 7 conseils que tout designer doit connaître sur l'accessibilité \url{https://medium.com/salesforce-ux/7-things-every-designer-needs-to-know-about-accessibility-64f105f0881b}
\item Pourquoi prioriser l'accessibilité web \url{https://blog.hubspot.com/service/why-prioritize-web-accessibility}
\end{itemize}

\textbf{Critères :}

\begin{itemize}
\item Exemples de mesures
\begin{itemize}
\item Contraste des couleurs
\item Texte mis en évidence ou agrandi au survol
\item Couleur des liens différente de celle du corps de texte
\item Un fil d'Ariane pour indiquer où vous êtes dans la navigation
\item Dans les formulaires, le titre du champ doit être au-dessus de la zone de texte
\item États d'erreur
\item etc.
\end{itemize}
\end{itemize}


\chapter{Exercice 03 : UI III - Ex 4 - Accessibilité (application)}
\exnumber{03}
\extitle{UI III - Ex 4 - Accessibilité (application)}
\exfiles{ex03/fichier Figma, nomme Name\_FirstName\_DeliverableName\_Date\_VersionNumber}
\exauthorize{Those named in the exercise}
\makeheaderfiles

\textbf{Instructions :}

Choose an app you find particularly accessible. Find at least 3 concrete examples of accessible measures and explain how they ease the reading or perception of visual content.

\textbf{Livrables :}

\begin{itemize}
\item A 3-page Figma file with 3 different elements from an app (it can be from 3 different apps) that are accessible. The three measures must be different.
\item Déposez le fichier Figma sur Github, nommé ainsi : « Name\_FirstName\_DeliverableName\_Date\_VersionNumber ». Exemple : « Doe\_John\_MidFiPrototype\_05242024\_V1 ».
\item Ajoutez un titre dans le fichier.
\item Soignez l'aspect global du fichier. Il doit être clair, simple et lisible. Vous pouvez regarder canva.com, sans surcharger le dessin de détails.
\end{itemize}

\textbf{Ressources :}

\begin{itemize}
\item What accessibility is and why it's so important \url{https://uxdesign.cc/what-accessibility-is-and-why-its-so-important-9c56e033ff26}
\item Web Content Accessibility Guidelines (WCAG) \url{https://www.w3.org/WAI/standards-guidelines/wcag/}
\item Accessibility testing  ---  W3C Wik \url{https://www.w3.org/wiki/Accessibility_testing}i
\item UX Myth: Accessibility is expensive and difficult \url{https://uxmyths.com/post/654091803/myth-5-accessibility-is-expensive-and-difficult}
\item Salesforce UX - 7 tips every designer needs to know about accessibility \url{https://medium.com/salesforce-ux/7-things-every-designer-needs-to-know-about-accessibility-64f105f0881b}
\end{itemize}

\textbf{Critères :}

\begin{itemize}
\item Examples of measures
\begin{itemize}
\item Color contrast
\item Highlighted or enlarged text when hovering
\item Links color that is different from the body text's
\item Breadcrumbs to indicate where you are in terms of navigation
\item In forms, the title of the case should be above the text zone
\item Error states
\item etc.
\end{itemize}
\end{itemize}


\chapter{Exercice 04 : UI III - Ex 5 - Fil d'Ariane}
\exnumber{04}
\extitle{UI III - Ex 5 - Fil d'Ariane}
\exfiles{ex04/fichier Figma, nomme Name\_FirstName\_DeliverableName\_Date\_VersionNumber}
\exauthorize{Those named in the exercise}
\makeheaderfiles

\textbf{Instructions :}

Browse the internet and find at least 5 different examples of breadcrumbs from websites or apps.

\textbf{Livrable :}

\begin{itemize}
\item A Figma file gathering 5 screenshots of breadcrumbs.
\item Déposez le fichier Figma sur Github, nommé ainsi : « Name\_FirstName\_DeliverableName\_Date\_VersionNumber ». Exemple : « Doe\_John\_MidFiPrototype\_05242024\_V1 ».
\item Ajoutez un titre dans le fichier.
\item Soignez l'aspect global du fichier. Il doit être clair, simple et lisible. Vous pouvez regarder canva.com, sans surcharger le dessin de détails.
\end{itemize}

\textbf{Ressources :}

\begin{itemize}
\item What are breadcrumbs? \url{https://www.seoptimer.com/blog/breadcrumbs-website/}
\end{itemize}


\chapter{Exercice 05 : UI III - Ex 6 - Interrupteurs}
\exnumber{05}
\extitle{UI III - Ex 6 - Interrupteurs}
\exfiles{ex05/fichier Figma, nomme Name\_FirstName\_DeliverableName\_Date\_VersionNumber}
\exauthorize{Those named in the exercise}
\makeheaderfiles

\textbf{Instructions :}

Browse the internet and find at least 5 different examples of toggle buttons from websites or apps.

\textbf{Livrable :}

\begin{itemize}
\item A Figma file gathering 5 screenshots of toggle buttons.
\item Déposez le fichier Figma sur Github, nommé ainsi : « Name\_FirstName\_DeliverableName\_Date\_VersionNumber ». Exemple : « Doe\_John\_MidFiPrototype\_05242024\_V1 ».
\item Ajoutez un titre dans le fichier.
\item Soignez l'aspect global du fichier. Il doit être clair, simple et lisible. Vous pouvez regarder canva.com, sans surcharger le dessin de détails.
\end{itemize}

\textbf{Ressources :}

\begin{itemize}
\item Toggle Switch - Tips \url{https://uxplanet.org/toggle-switch-5-simple-design-tips-for-better-design-b4046eff4a2f}
\end{itemize}


\chapter{Exercice 06 : UI III - Ex 7 - Boutons radio}
\exnumber{06}
\extitle{UI III - Ex 7 - Boutons radio}
\exfiles{ex06/fichier Figma, nomme Name\_FirstName\_DeliverableName\_Date\_VersionNumber}
\exauthorize{Those named in the exercise}
\makeheaderfiles

\textbf{Instructions :}

Browse the internet and find at least 5 different examples of radio buttons from websites or apps.

\textbf{Livrable :}

\begin{itemize}
\item A Figma file gathering 5 screenshots of radio buttons.
\item Déposez le fichier Figma sur Github, nommé ainsi : « Name\_FirstName\_DeliverableName\_Date\_VersionNumber ». Exemple : « Doe\_John\_MidFiPrototype\_05242024\_V1 ».
\item Ajoutez un titre dans le fichier.
\item Soignez l'aspect global du fichier. Il doit être clair, simple et lisible. Vous pouvez regarder canva.com, sans surcharger le dessin de détails.
\end{itemize}

\textbf{Ressources :}

\begin{itemize}
\item What is a radio button? \url{https://www.justinmind.com/blog/radio-button-design-examples/}
\end{itemize}


\chapter{Exercice 07 : UI III - Ex 8 - Calendriers}
\exnumber{07}
\extitle{UI III - Ex 8 - Calendriers}
\exfiles{ex07/fichier Figma, nomme Name\_FirstName\_DeliverableName\_Date\_VersionNumber}
\exauthorize{Those named in the exercise}
\makeheaderfiles

\textbf{Instructions :}

Browse the internet and find at least 5 different examples of calendars from websites or apps.

\textbf{Livrable :}

\begin{itemize}
\item A Figma file gathering 5 screenshots of calendars.
\item Déposez le fichier Figma sur Github, nommé ainsi : « Name\_FirstName\_DeliverableName\_Date\_VersionNumber ». Exemple : « Doe\_John\_MidFiPrototype\_05242024\_V1 ».
\item Ajoutez un titre dans le fichier.
\item Soignez l'aspect global du fichier. Il doit être clair, simple et lisible. Vous pouvez regarder canva.com, sans surcharger le dessin de détails.
\end{itemize}


\chapter{Exercice 08 : UI III - Ex 9 - Sélecteurs d'heure}
\exnumber{08}
\extitle{UI III - Ex 9 - Sélecteurs d'heure}
\exfiles{ex08/fichier Figma, nomme Name\_FirstName\_DeliverableName\_Date\_VersionNumber}
\exauthorize{Those named in the exercise}
\makeheaderfiles

\textbf{Instructions :}

Browse the internet and find at least 5 different examples of time pickers from websites or apps.

\textbf{Livrable :}

\begin{itemize}
\item A Figma file gathering 5 screenshots of time pickers.
\item Déposez le fichier Figma sur Github, nommé ainsi : « Name\_FirstName\_DeliverableName\_Date\_VersionNumber ». Exemple : « Doe\_John\_MidFiPrototype\_05242024\_V1 ».
\item Ajoutez un titre dans le fichier.
\item Soignez l'aspect global du fichier. Il doit être clair, simple et lisible. Vous pouvez regarder canva.com, sans surcharger le dessin de détails.
\end{itemize}


\chapter{Exercice 09 : UI III - Ex 10 - Micro-interactions}
\exnumber{09}
\extitle{UI III - Ex 10 - Micro-interactions}
\exfiles{ex09/fichier Figma, nomme Name\_FirstName\_DeliverableName\_Date\_VersionNumber}
\exauthorize{Those named in the exercise}
\makeheaderfiles

\textbf{Instructions :}

Browse the internet and find at least 5 different examples of feedback micro-interactions or micro-animations on websites or apps.

Be careful! Feedback is helpful to provide information to the user, not just showing a catchy design. You need to make the difference.

\textbf{Livrable :}

\begin{itemize}
\item A Figma file gathering 5 examples of \textbf{feedback} micro-interactions or animations. In each example, show the different versions of the element.
\item Déposez le fichier Figma sur Github, nommé ainsi : « Name\_FirstName\_DeliverableName\_Date\_VersionNumber ». Exemple : « Doe\_John\_MidFiPrototype\_05242024\_V1 ».
\item Ajoutez un titre dans le fichier.
\item Soignez l'aspect global du fichier. Il doit être clair, simple et lisible. Vous pouvez regarder canva.com, sans surcharger le dessin de détails.
\end{itemize}

\textbf{Ressources :}

\begin{itemize}
\item Why use micro-interactions in your design \url{https://hike.one/update/why-use-micro-animations-in-your-design}
\item Micro-interactions: why, when,  and how to use them to boost the UX \url{https://uxdesign.cc/micro-interactions-why-when-and-how-to-use-them-to-boost-the-ux-17094b3baaa0}
\end{itemize}

\textbf{Pistes :}

\begin{itemize}
\item Feedback is the situation when the app specifically sends a piece of information to the user, like: 
\item Payment accepted
\item Refreshed feed
\item Information saved
\item etc.
\end{itemize}


\chapter{Exercice 10 : UI III - Ex 11 - Cohérence}
\exnumber{10}
\extitle{UI III - Ex 11 - Cohérence}
\exfiles{ex10/fichier Figma, nomme Name\_FirstName\_DeliverableName\_Date\_VersionNumber}
\exauthorize{Those named in the exercise}
\makeheaderfiles

\textbf{Instructions :}

Pick a website or an app of your choice*. Highlight 5 signs of consistency across the pages.

*Airbnb, Google, Apple, Amazon, Facebook, Instagram and WhatsApp are not accepted.

\textbf{Livrable :}

\begin{itemize}
\item A Figma file gathering at least 5 signs of consistency.
\item Déposez le fichier Figma sur Github, nommé ainsi : « Name\_FirstName\_DeliverableName\_Date\_VersionNumber ». Exemple : « Doe\_John\_MidFiPrototype\_05242024\_V1 ».
\item Ajoutez un titre dans le fichier.
\item Soignez l'aspect global du fichier. Il doit être clair, simple et lisible. Vous pouvez regarder canva.com, sans surcharger le dessin de détails.
\end{itemize}

\textbf{Ressources :}

\begin{itemize}
\item 7 tips for designing consistency \url{https://designshack.net/articles/graphics/7-tips-for-designing-consistency/}
\item The value of consistent design \url{https://www.invisionapp.com/inside-design/consistent-design/}
\item Consistency and Standards - Examples of Nielsens's Design Heuristics \url{https://medium.com/@gregoralbrecht/consistency-and-standards-nielsens-design-heuristic-explained-ac91f450fd8}
\item The importance of consistency in design work \url{https://yesimadesigner.com/the-importance-of-consistancy-in-design-work/}
\item Examples of signs of consistency: 
\begin{itemize}
\item colors
\item same shaped buttons
\item alignment
\item consistent product pages (with identical frameworks)
\item etc.
\end{itemize}
\end{itemize}


\end{document}
